\subsection{Index construction: weights, divisors, returns and index types}

These entries explain how an index level is computed and maintained, and which kinds of equity and bond index exist. The worked example uses three invented stocks A, B and C and two invented bonds; its arithmetic is checked in \texttt{trading\_notes.numerical.index\_construction.index\_construction\_checks}. Major regional indices and the instruments that track them are in \hyperref[def:equity-indices-derivatives]{equity indices and instruments}.

\begin{qfmethod}[Index values and returns: price, capitalisation and equal weighting, the divisor and its adjustments, corporate actions, price, total and net total return, rebalancing and turnover]\label{meth:index-value-return-calculation}
\qffield{Statement} An index level is a market value scaled by a divisor. For a capitalisation-weighted index S\&P Dow Jones Indices write
\[ \text{Index level}=\frac{\sum_i P_iQ_i}{\text{Divisor}},\qquad Q_i=\text{IWF}_i\times\text{total shares}_i, \]
where $P_i$ is the price, $Q_i$ the share count used in the index and the investable weight factor (IWF) the fraction of shares available to investors, so the index is float-adjusted \citep[pp.~6--7, eqs.~(1) and~(3)]{index-mathematics-methodology-2026}. The formula is a modified Laspeyres index whose divisor ``both represents the initial market value and sets the base value'' (p.~6). ``The purpose of the index divisor is to maintain the continuity of an index level following the implementation of corporate actions, index rebalancing events, or other non-market driven actions'' (p.~5).

When a change alters the index market value at unchanged prices (an addition, a deletion, a share issue, an IWF change), the divisor is reset after the close so that the level is the same before and after:
\[ \text{Divisor}_{t}=\text{Divisor}_{t-1}\frac{MV_t}{MV_{t-1}}\quad\text{or}\quad \text{Divisor}_{\text{new}}=\text{Divisor}_{\text{old}}+\frac{\text{CMV}}{\text{Index level}}, \]
with CMV the change in market value; the additive form lets separate changes be summed (pp.~7--8, eqs.~(6)--(7)).

\emph{Weighting schemes.} In a \emph{price-weighted} index, such as the Dow Jones Industrial Average, ``constituent weights are determined solely by the prices of the constituent stocks'', and the divisor adjusts ``for any price impacting corporate action'', including stock splits (p.~15). An \emph{equal-weighted} index gives every stock the same weight at each rebalancing through an additional weight factor $\text{AWF}_i=Z/(N\cdot\text{FAMV}_i)$; ``As stock prices move, the weights will shift, and exact equality will be lost'', so it ``must be rebalanced from time to time'', and the divisor changes at the rebalancing (p.~16, eqs.~(2)--(3)). Capped and other non-capitalisation weightings use the same AWF device with user-defined target weights (pp.~10--14).

\emph{Corporate actions.} A stock split multiplies shares and divides the price by the same factor, ``such that the market capitalization of the company remains the same'' \citep[p.~24]{equity-index-policies-2026}, so a capitalisation-weighted divisor does not change. An ordinary dividend ``will only have an effect on the Total Return (TR) and Net Total Return (NTR) indices and not on Price Return indices''; a special dividend is treated as a corporate action with ``price and divisor adjustments'' (p.~19). A spin-off is generally added at a zero price at the close before the ex-date ``with no divisor adjustment'' (p.~8).

\emph{Total return.} The dividends going ex on a day are converted into index points, $\text{IndexDividend}=\sum_i D_iQ_i/\text{Divisor}$, and the daily total return is $\text{DTR}_t=(\text{Level}_t+\text{IndexDividend}_t)/\text{Level}_{t-1}-1$; equivalently $\text{TR}_t=\text{TR}_{t-1}(\text{PR}_t+\text{IndexDividend}_t)/\text{PR}_{t-1}$. The total return index reinvests dividends ``in the overall index, not in the specific stock paying the dividend''. A net total return index deducts withholding tax from each dividend before the same steps \citep[pp.~21--22]{index-mathematics-methodology-2026}. S\&P DJI's global indices use the maximum non-resident withholding rates without treaty relief \citep[p.~24]{equity-index-policies-2026}.

\emph{Turnover.} One-way index turnover is $\tfrac12\sum_i|w_i^{\text{close}}-w_i^{\text{adjusted}}|$ over constituents, excluding weight changes caused by price moves \citep[p.~70]{index-mathematics-methodology-2026}.

\qffield{Notation} $MV$ index market value; $w_i=P_iQ_i/MV$ capitalisation weight; with no changes between two closes, the index return equals $\sum_iw_{i,t-1}r_{i,t}$ for a capitalisation- or equal-weighted index, with weights at the previous close.

\qffield{Practice} Example, original (base values 100 and 1{,}000; prices in currency units, shares in millions). Day~0 prices are A 50, B 20, C 80; day~1 prices 55, 19, 84; shares 100, 300, 50; IWFs 0.8, 1.0, 0.5.
\begin{itemize}
\item \emph{Price-weighted.} The divisor is $150/100=1.5$. Day~1: $158/1.5=105.333$, up 5.333 percent. C carries 53.3 percent of the weight only because its price is highest. After the close C splits 2-for-1 (price 84 to 42). The new divisor $116/105.333=1.10127$ keeps the level at 105.333. On day~2 (54, 20, 44) the level is $118/1.10127=107.149$, up 1.724 percent. Without the adjustment it would print 78.67, a fall of 25.3 percent that no investor suffered, and C's weight drops from 53.2 to 36.2 percent at the split.
\item \emph{Capitalisation-weighted.} Float-adjusted market value is $50\times80+20\times300+80\times25=12{,}000$, so the divisor is 12. Day~1: $12{,}200/12=1{,}016.67$, up 1.667 percent. This equals the weighted return $\tfrac13(10\%)+\tfrac12(-5\%)+\tfrac16(5\%)$. After the close, C's split changes nothing, but B issues 30 million shares, raising market value by $19\times30=570$. The divisor becomes $12+570/1{,}016.67=12.5607$. Day~2 market value is 13{,}120, so the level is 1{,}044.53, up 2.741 percent.
\item \emph{Total and net total return, day~2.} B goes ex-dividend 0.40. The dividend points are $0.40\times330/12.5607=10.509$. The total return is $(1{,}044.53+10.51)/1{,}016.67-1=3.774$ percent. With 15 percent withholding the net total return is 3.619 percent, against a price return of 2.741 percent.
\item \emph{Equal-weighted.} The day~1 return is the simple average, 3.333 percent. Weights drift to 35.48, 30.65 and 33.87 percent. Rebalancing back to one third each needs one-way turnover of 2.69 percent before day~2.
\end{itemize}

\qffield{Conditions} The formulas are S\&P DJI's for its indices as of September 2026; other providers' share-count, float and timing rules differ, and an index's own methodology governs its corporate-action treatment. The level is unchanged by a maintenance event only at the prices used for the adjustment, normally the closing prices. Quarterly share updates, buffers and capping are index-specific.

\qffield{Pitfalls} Reading a price index return as an investor's return omits dividends; comparing a total return portfolio with a price index flatters the portfolio. Price weighting makes weights depend on share-price levels, so a split cuts a stock's weight. Equal and capped weightings need rebalancing trades that a capitalisation-weighted index avoids, and those trades cost a tracking fund money that the index does not bear. A total return index assumes reinvestment across the whole index on the ex-date, not in the paying stock on the pay date.

\qffield{See also} \hyperref[def:equity-fixed-income-index-types]{index types and uses}; \hyperref[def:return-measures]{return measures} (price and total return); \hyperref[def:equity-indices-derivatives]{equity indices and instruments}.

\end{qfmethod}

\begin{qfdefinition}[Index types and uses: broad, size, style, sector, factor, equal-weighted, fundamental and optimised equity indices; aggregate, government, corporate, high-yield and inflation-linked bond indices; benchmarks, asset-class proxies and the basis of index products]\label{def:equity-fixed-income-index-types}
\qffield{Statement} Bacon lists the calculation families of commercial equity indices. A \emph{price-weighted} (aggregate price) index adds prices and divides by a divisor. A \emph{geometric} index takes the $n$th root of the product of price relatives and cannot be held in a real portfolio. In a \emph{market capitalisation} index, usually free-float adjusted and computed as a Laspeyres index, each security's weight is set by its value. An \emph{equal-weighted} index needs periodic rebalancing whose frequency ``must also be arbitrarily chosen''. A \emph{fundamental} index weights by book value, sales, cash flow, employees or dividends. An \emph{optimised} index (minimum variance, efficient) embeds part of an investment process. \emph{Style- and factor-based} indices split the market by size and by value and growth \citep[3rd ed., 2023, chap.~4, pp.~105--111]{wide-bacon-attribution}. S\&P DJI separate capitalisation-weighted indices from non-capitalisation ones, which use ``equal weighting, factor weighting such as dividend yield or volatility, strategic tilts, thematic weighting, or other alternative weighting schemes'', and from derived indices such as total return, leveraged, inverse and risk-control indices \citep[p.~4]{index-mathematics-methodology-2026}.

\emph{Fixed-income indices} are categorised by issuer (government, corporate, municipal, asset-backed and mortgage-backed), country, currency, credit rating, liquidity, maturity and sector, and report coupon, life, duration, convexity and yield as well as return \citep[pp.~111--112]{wide-bacon-attribution}. Bloomberg's flagship indices are rules-based baskets designed to be ``Representative'', ``Replicable'', ``Objective and transparent'' and ``Relevant'', weighted by market value of amount outstanding and rebalanced monthly. A fixed ``Returns Universe'' set at the start of the month determines that month's return; a daily ``Projected Universe'' previews the next \citep[pp.~9 and~42--43]{fixed-income-index-methodology-2026}. A bond's monthly total return is price return plus coupon return (plus paydown return for amortising bonds), each on the beginning dirty value, with clean prices in the price return. The index return is $\sum(\text{bond return}\times\text{beginning market-value weight})$, and intra-month cash earns no reinvestment return until the month-end reset (pp.~47 and~57--61).

\qffield{Practice} \emph{Uses.} Bloomberg observes three uses of bond indices: ``Portfolio performance targets'', ``Informational measures of asset class risk and return'' and ``References for index-linked products'' \citep[p.~9]{fixed-income-index-methodology-2026}. Equity indices serve the same roles. A broad float-adjusted index proxies an equity market in asset allocation and is the benchmark of a core mandate. Style, size and sector indices benchmark specialists and are the regressors of returns-based style analysis (\hyperref[meth:manager-due-diligence]{manager due diligence}). Indices also underlie index funds, ETFs, futures and options (\hyperref[def:equity-indices-derivatives]{equity indices and instruments}). Which type fits depends on the question: a capitalisation-weighted index measures the market's aggregate experience; an equal-weighted or factor index measures a strategy.

Example, original (two-bond month). Bond 1 has par 500, clean price 98.00 to 99.00 and accrued interest 1.00 to 1.40. Its price return is $1.00/99.00=1.010$ percent and its coupon return $0.40/99.00=0.404$ percent, a total of 1.414 percent. Bond 2 has par 300, clean price 102.50 to 101.80 and accrued 0.50 to 0.90, for a total return of $-0.291$ percent. Beginning market values are 495 and 309, weights 61.6 and 38.4 percent, and the index returns 0.759 percent for the month. A bond issued mid-month joins only at the next reset.

\qffield{Conditions} Bacon's judgements, such as that price-only indices ``should never be used'' for performance analysis and that equal-weighted and fundamental indices ``tend to outperform'' because they avoid overweighting overpriced stocks, are his views, not results proved in the text (pp.~106 and~110). Bloomberg's rules are as of the January 2026 version; its two-business-day lockout before month end began with the August 2025 Returns Universe (p.~42).

\qffield{Pitfalls} Debt-weighted bond indices give the largest weights to the most indebted issuers, the ``bums'' problem, and their duration reflects issuers' choices, not investors' liabilities \citep[p.~112]{wide-bacon-attribution}. Many bonds trade rarely, so index prices are partly evaluated, not traded. Two providers' indices with the same label can differ in sampling, filtering, buffering and free-float rules (p.~113). An optimised or factor index already contains investment decisions, so a manager measured against it is measured against a strategy.

\qffield{See also} \hyperref[meth:index-value-return-calculation]{index values and returns}; \hyperref[meth:benchmark-selection-quality]{benchmark selection and quality}; \hyperref[def:commodity-indices-risk-return]{commodity indices}.

\end{qfdefinition}
